\subsection{\texorpdfstring{第~\ref{sec:cosmology}~节}{宇宙学}}\label{subsec:solutions-cosmology}

\begin{description}[font=\normalfont,leftmargin=2em,labelwidth=1.5em]
\RGsolution{ex:hubble-law-and-cosmological-redshift}{hubble-law-and-cosmological-redshift}
\begin{enumerate}[label=(\alph*),ref=\alph*,leftmargin=*]
\item 两个星系之间的距离显然为 $d(t)=a(t)x$, 其中
\[
  x=\int_0^{r_1}\frac{\mathrm{d}r}{\sqrt{1-kr^2}}.
\]
因此
\[
  \dot d=\dot a x=\frac{\dot a}{a}ax=Hd.
\]
\item 这些曲线的切向量为
\[
  \dot c=\dot t\frac{\partial}{\partial t}+\frac{\partial}{\partial r},
\]
它是零模向量当且仅当
\[
  -\dot t^2+\frac{a^2(t)}{1-kr^2}=0
  \quad\Longleftrightarrow\quad \dot t=\pm\frac{a(t)}{\sqrt{1-kr^2}}.
\]
对于从第一个星系连到第二个星系的零模测地线, 必须取正号. 因此
\[
  \dot c=\frac{a(t)}{\sqrt{1-kr^2}}\frac{\partial}{\partial t}
    +\frac{\partial}{\partial r}=\frac{a(t)}{\sqrt{1-kr^2}}(E_0+E_r).
\]
由于
\[
  \omega_r^0=\omega_0^r=\frac{\dot a}{a}\omega^r,
\]
可知
\begin{align*}
  \nabla_{E_0+E_r}(E_0+E_r)
    &=\omega_0^r(E_0+E_r)E_r+\omega_r^0(E_0+E_r)E_0\\
    &=\frac{\dot a}{a}(E_0+E_r).
\end{align*}
因此 $E_0+E_r$ 与一条重新参数化的测地线相切, 从而 $\dot c$ 也与之相切.
\item 对 $\frac{\mathrm{d}t}{\mathrm{d}r}$ 所满足的方程求导, 得
\[
  \frac{\mathrm{d}t'}{\mathrm{d}r}
     =\frac{\dot a(t)t'}{\sqrt{1-kr^2}}
     =\frac{\dot a(t)}{a(t)}\frac{\mathrm{d}t}{\mathrm{d}r}t',
\]
其中 $t'=\frac{\partial t}{\partial t_0}$. 积分得
\[
  \int_0^{r_1}\frac{1}{t'}\frac{\mathrm{d}t'}{\mathrm{d}r}\,\mathrm{d}r
    =\int_0^{r_1}\frac{\dot a(t)}{a(t)}\frac{\mathrm{d}t}{\mathrm{d}r}\,\mathrm{d}r
    \quad\Longleftrightarrow\quad
    \log(t'(r_1))=\log\left(\frac{a(t_1)}{a(t_0)}\right).
\]
\end{enumerate}

\RGsolution{ex:einstein-and-de-sitter-universe-embeddings}{einstein-and-de-sitter-universe-embeddings}
在球坐标下, 四维 Minkowski 时空的度量可写为
\[
  g=-\mathrm{d}t\otimes\mathrm{d}t+\mathrm{d}r\otimes\mathrm{d}r+r^2h,
\]
其中 $h$ 是 $S^3$ 上的标准度量.
\begin{enumerate}[label=(\alph*),ref=\alph*,leftmargin=*]
\item “柱面”的方程为 $r^2=\frac{1}{\Lambda}$, 因而诱导度量为
\[
  g_{Einstein}=-\mathrm{d}t\otimes\mathrm{d}t+\frac{1}{\Lambda}h.
\]
这就是 $k=1$ 且 $a(t)=\frac{1}{\sqrt{\Lambda}}$ 的 FLRW 模型的度量. 若 $\alpha=\frac{1}{3\sqrt{\Lambda}}$ [参见 Exercise~\ref{ex:cosmology}(\ref{ex:flrw-with-cosmological-constant})], 则在这种情形下,
\[
  \frac{\dot a^2}{2}-\frac{\alpha}{a}-\frac{\Lambda}{6}a^2
    =-\alpha\sqrt{\Lambda}-\frac{1}{6}=-\frac{1}{2}.
\]
因此, 这个度量满足密度为
\[
  \rho=\frac{3\alpha}{4\pi a^3}=\frac{\Lambda}{4\pi}
\]
的无压流体所对应的 Einstein 方程, 其宇宙学常数 $\Lambda>0$.
\item “球面”的方程为 $-t^2+r^2=\frac{3}{\Lambda}$, 其解为
\[
  \begin{cases}
    t=\sqrt{\dfrac{3}{\Lambda}}\sinh\left(\sqrt{\dfrac{\Lambda}{3}}\,\tau\right)\\[.8ex]
    r=\sqrt{\dfrac{3}{\Lambda}}\cosh\left(\sqrt{\dfrac{\Lambda}{3}}\,\tau\right).
  \end{cases}
\]
因此, 诱导度量为
\[
  g_{deSitter}=-\mathrm{d}\tau\otimes\mathrm{d}\tau
     +\frac{3}{\Lambda}\cosh^2\left(\sqrt{\frac{\Lambda}{3}}\,\tau\right)h.
\]
这就是 $k=1$ 且 $a(\tau)=\sqrt{\frac{3}{\Lambda}}\cosh\left(\sqrt{\frac{\Lambda}{3}}\,\tau\right)$ 的 FLRW 模型的度量. 若 $\alpha=0$, 则在这种情形下,
\[
  \frac{\dot a^2}{2}-\frac{\alpha}{a}-\frac{\Lambda}{6}a^2
    =-\frac{1}{2}-\frac{\alpha}{a}=-\frac{1}{2}.
\]
因此, 这个度量满足密度 $\rho=0$ 的无压流体所对应的 Einstein 方程, 其宇宙学常数 $\Lambda>0$.
\end{enumerate}

\RGsolution{ex:flrw-brightness-and-redshift}{flrw-brightness-and-redshift}
\begin{enumerate}[label=(\alph*),ref=\alph*,leftmargin=*]
\item 由 Exercise~\ref{ex:cosmology}(\ref{ex:hubble-law-and-cosmological-redshift}) 可知, 连接两个星系的零模测地线可通过求解微分方程
\[
  \frac{\mathrm{d}t}{\mathrm{d}r}=\frac{t}{\sqrt{1+r^2}}
\]
找到. 这个方程可分离变量, 积分得到
\begin{align*}
  &\int_{t_0}^{t_1}\frac{\mathrm{d}t}{t}=\int_0^{r_1}\frac{\mathrm{d}r}{\sqrt{1+r^2}}\\
  &\quad\Longleftrightarrow\quad
    \log\left(\frac{t_1}{t_0}\right)=\operatorname{arcsinh}(r_1)\\
  &\quad\Longleftrightarrow\quad
    r_1=\sinh(\log(1+z))=\frac{1}{2}\left(1+z-\frac{1}{1+z}\right).
\end{align*}
因此
\[
  R=t_1r_1=t_1\frac{z^2+2z}{2+2z}.
\]
\item 在这种情形下, 必须求解
\[
  \frac{\mathrm{d}t}{\mathrm{d}r}=e^{Ht},
\]
于是
\[
  \int_{t_0}^{t_1}e^{-Ht}\,\mathrm{d}t=\int_0^{r_1}\mathrm{d}r
  \quad\Longleftrightarrow\quad
  r_1=\frac{1}{H}\bigl(e^{-Ht_0}-e^{-Ht_1}\bigr).
\]
因此
\[
  R=e^{Ht_1}r_1=\frac{z}{H}.
\]
\item 在这种情形下, 必须求解
\[
  \frac{\mathrm{d}t}{\mathrm{d}r}=\left(\frac{t}{t_1}\right)^{\frac{2}{3}},
\]
于是
\begin{align*}
  &\int_{t_0}^{t_1}\left(\frac{t}{t_1}\right)^{-\frac{2}{3}}\,\mathrm{d}t=\int_0^{r_1}\mathrm{d}r
    \quad\Longleftrightarrow\\
  &\qquad r_1=3t_1^{\frac{2}{3}}\bigl(t_1^{\frac{1}{3}}-t_0^{\frac{1}{3}}\bigr)
     =3t_1\left(1-(1+z)^{-\frac{1}{2}}\right).
\end{align*}
因此
\[
  R=r_1=3t_1\left(1-(1+z)^{-\frac{1}{2}}\right).
\]
\end{enumerate}
\end{description}
